% U3 3절 — 라플라스 변환의 최소 문법 (정본: paper/sections/03_lti_system.tex 발췌 재구성)
\section{라플라스 변환의 최소 문법}
\label{sec:u3-laplace}

앞 절의 전신선로가 남긴 질문 --- 신호의 크기와 위상이 왜, 어떻게 변하는가 ---
을 계산 가능한 형태로 바꾸려면 두 가지 준비가 필요하다.
하나는 시스템에 대한 가정이고, 다른 하나는 미분방정식을 대수식으로 바꾸는
변환이다.
이 절은 그 최소한의 문법을 유도 생략 없이 정리한다.

먼저 가정부터 보자.
시스템이 선형(linear)이라는 것은, 입력 $u(t)$를 출력 $y(t)$로 바꾸는 규칙
$\mathcal{S}$가 임의의 입력 $u_1$, $u_2$와 스칼라 $a$, $b$에 대해
$\mathcal{S}\{a u_1 + b u_2\} = a\,\mathcal{S}\{u_1\} + b\,\mathcal{S}\{u_2\}$를
만족한다는 뜻이다.
입력을 두 배로 하면 출력도 두 배가 되고, 두 입력을 더해 넣은 결과는 각각의
출력을 더한 것과 같다.
시불변(time-invariant)이라는 것은 같은 입력을 $t_0$만큼 늦게 넣으면 출력도
똑같이 $t_0$만큼 늦어질 뿐 모양은 변하지 않는다는 뜻,
즉 $\mathcal{S}\{u(t-t_0)\} = y(t-t_0)$이다.
실제 로봇은 완전한 선형도 시불변도 아니지만, 작은 움직임과 한 작동점
근처에서는 이 선형시불변(LTI) 가정이 많은 현상을 충분히 잘 설명하며,
전달함수와 주파수 응답이라는 강력한 도구를 쓸 수 있게 해 준다.

\subsection{정의에서 미분 규칙까지}

라플라스 변환은 시간 함수 $x(t)$를 복소 변수 $s$의 함수로 요약하는 적분이다.
\begin{equation}
  X(s) = \mathcal{L}\{x(t)\}
  = \int_{0}^{\infty} x(t)\,e^{-st}\,\dd t.
  \label{eq:u3-laplace-def}
\end{equation}
말로 풀면, $x(t)$에 $e^{-st}$라는 점점 작아지는 추를 곱해 $t=0$부터 끝까지
모두 더한 값이다.
$s$를 하나 고를 때마다 숫자가 하나 나오므로, 시간 파형 전체가 $X(s)$ 한
장으로 요약된다.

이 변환이 유용한 이유는 미분이 $s$를 곱하는 연산으로 바뀌기 때문인데,
이 규칙은 외울 공식이 아니라 정의 \eqref{eq:u3-laplace-def}와 곱의 미분
규칙만으로 나온다.
$s$는 $t$와 무관하므로
\begin{equation}
  \frac{\dd}{\dd t}\left[x(t)\,e^{-st}\right]
  = \dot{x}(t)\,e^{-st} - s\,x(t)\,e^{-st}
\end{equation}
이다.
양변을 $t=0$부터 $\infty$까지 적분하자.
왼쪽은 미분한 것을 다시 적분하므로 끝값 빼기 시작값,
즉 $x(t)e^{-st}$의 $t\to\infty$ 값에서 $t=0$ 값 $x(0)$을 뺀 것이다.
응답이 발산하지 않아 $t\to\infty$에서 $x(t)e^{-st}\to 0$이 되는 $s$에서 보면
왼쪽은 $0 - x(0) = -x(0)$이다.
오른쪽의 두 적분은 정의에 따라 각각 $\mathcal{L}\{\dot{x}\}$와 $-sX(s)$이므로
$-x(0) = \mathcal{L}\{\dot{x}\} - sX(s)$, 정리하면
\begin{equation}
  \mathcal{L}\{\dot{x}(t)\} = sX(s) - x(0)
  \label{eq:u3-deriv-rule}
\end{equation}
이다.
초기조건을 0으로 두면 $\dot{x}$는 $sX(s)$처럼 다룰 수 있다.

간단한 검산을 해 보자.
$x(t)=1$(상수)이면 정의에서
$X(s)=\int_0^\infty e^{-st}\dd t
=\left[-e^{-st}/s\right]_0^\infty
=0-\left(-\tfrac{1}{s}\right)=\tfrac{1}{s}$이다.
$\dot{x}=0$이므로 규칙 \eqref{eq:u3-deriv-rule}의 왼쪽은
$\mathcal{L}\{0\}=0$이고, 오른쪽도
$sX(s)-x(0)=s\cdot(1/s)-1=0$으로 일치한다.

두 번 미분한 항은 같은 규칙을 두 번 적용하면 된다.
$\dot{x}$를 새 함수로 보고 \eqref{eq:u3-deriv-rule}를 다시 쓰면
\begin{equation}
  \mathcal{L}\{\ddot{x}\}
  = s\,\mathcal{L}\{\dot{x}\} - \dot{x}(0)
  = s^2 X(s) - s\,x(0) - \dot{x}(0)
\end{equation}
이고, 초기조건을 0으로 두면 $\ddot{x}$는 $s^2 X(s)$처럼 다룰 수 있다.

\subsection{미분방정식에서 전달함수로}

이제 질량-스프링-댐퍼의 운동방정식 $m\ddot{x}+b\dot{x}+kx=f$의 각 항을
초기조건 0에서 하나씩 바꾸면
$m\ddot{x}\to m s^2 X(s)$, $b\dot{x}\to b s X(s)$, $kx\to kX(s)$,
$f\to F(s)$이므로
\begin{equation}
  (m s^2 + b s + k)\,X(s) = F(s)
\end{equation}
가 된다.
시간에 따라 계속 갱신해야 하는 미분 장부가 한 장의 대수 장부로 바뀐 것이다.
입력과 출력의 비율로 정리하면 전달함수(transfer function)
\begin{equation}
  H(s) = \frac{X(s)}{F(s)} = \frac{1}{m s^2 + b s + k}
  \label{eq:u3-msd-tf}
\end{equation}
를 얻는다.
분모의 $ms^2$ 항은 관성, $bs$ 항은 감쇠, $k$ 항은 스프링 복원의 영향을 담는다.
전달함수는 초기 저장 에너지를 0으로 둔(zero-state) 입력-출력 관계이며,
같은 물리 장치라도 무엇을 입력과 출력으로 고르느냐에 따라 다른 전달함수로
보인다는 점만 기억해 두자.

여기서 $s$는 단순한 알파벳이 아니라 시간응답의 성질을 담는 복소 변수로,
보통 $s=\sigma+\jj\omega$로 생각한다.
$\sigma$는 응답이 커지거나 줄어드는 속도, $\omega$는 흔들리는 속도와
관련된다.
실제 물리 입력이 복소수라는 뜻이 아니라, 감쇠와 진동을 한 기호 안에서 함께
계산하기 위한 표현이다.

\subsection{주파수 응답 읽기}

안정한 LTI 시스템에 정현파 입력 $u(t)=A\cos(\omega t)$를 넣으면, 정상상태
출력은 같은 주파수를 가지며 크기와 위상만 달라진다.
\begin{equation}
  y_{\mathrm{ss}}(t)
  = A\,|H(\jj\omega)|\cos\!\left(\omega t + \angle H(\jj\omega)\right).
\end{equation}
$|H(\jj\omega)|$와 $\angle H(\jj\omega)$는 복소수 $H(\jj\omega)$의 크기와
각도이다.
복소수 $a+\jj b$를 평면 위의 점 $(a,b)$로 그리면 크기는 원점까지의 거리
$\sqrt{a^2+b^2}$, 각도는 양의 실수축에서 잰 회전각이다.
예를 들어 어떤 주파수에서 $H(\jj\omega)=1-\jj$라면 크기는
$\sqrt{1^2+(-1)^2}=\sqrt{2}$, 각도는 $-45^\circ$이므로, 출력 진폭은 입력의
약 1.41배가 되고 위상은 $45^\circ$만큼 늦는다.
즉 $s=\jj\omega$를 대입한 $H(\jj\omega)$ 하나로, 각 주파수에서 시스템이
얼마나 증폭하거나 감쇠하는지, 얼마나 늦게 반응하는지를 모두 읽을 수 있다.
pole과 응답 모양의 관계, 컨볼루션, 최종값 정리 등 여기서 줄인 내용의 세부
유도는 확장판을 참고한다.

이제 이 문법을 첫 번째 실물에 적용할 차례다.
다음 절에서는 저항, 인덕터, 커패시터로 이루어진 RLC 회로의 방정식을 같은
절차로 변환하여, 전기 임피던스가 어떻게 읽히는지 본다.
